\section{A deployable controller}
\label{sec:controller}

The benchmark of Section~\ref{sec:benchmark} evaluates the reference field at
every epoch of a reference trajectory it already knows. This section removes
that access and measures what each removal costs. This is RQ3.

\paragraph{What ``without the reference field'' means, precisely.} The
distinction is not between having and not having the gravity model --- a
propagator carries the coefficients by definition, and pretending otherwise
would be artificial. It is between \emph{global, offline model metadata} and
\emph{trajectory-specific evaluations of it}:

\begin{center}\small
\begin{tabular}{@{}lll@{}}
  \toprule
  & Available to the controller & Withheld \\
  \midrule
  Field & degree variances $P_n$ (a table of \ph{n} numbers) &
    $\Delta\mathbf a_{\mathrm{ref}}(t)$ along the arc \\
  Trajectory & a cheap pilot arc it computes itself & the reference arc \\
  Sensitivity & $\Phi$ from low-degree variational equations & the exact $\Phi$ \\
  \bottomrule
\end{tabular}
\end{center}

The controller may know the spectrum; it may not be told what the omitted
acceleration is at the epochs it will fly through. That is the information
boundary the rest of this section works against, and it is the one a real
mission faces.

\begin{table}[!htbp]
  \centering\small
  \caption{What the benchmark uses, and what replaces it.}
  \label{tab:substitutions}
  \begin{tabular}{@{}p{0.26\textwidth}p{0.34\textwidth}p{0.32\textwidth}@{}}
    \toprule
    Benchmark uses & Replaced by & Cost, charged to $B_+$ \\
    \midrule
    Reference field, pointwise & leading omitted bands (direction) plus the
      measured degree-variance spectrum (amplitude) & one synthesis per
      look-ahead probe point; measured \ph{pct} --- the dominant term \\
    Reference arc & cheap pilot arc, $N=40$, loose tolerance & measured
      \ph{pct} of the run's gravity work \\
    Exact $\Phi$ & low-degree variational equations integrated on the pilot
      arc, $\dot\Phi = \bigl[\begin{smallmatrix}\mathbf 0&\mathbf I\\
      \mathbf G_{40}&\mathbf 0\end{smallmatrix}\bigr]\Phi$ & \ph{pct} \\
    \bottomrule
  \end{tabular}
\end{table}

\subsection{The omitted acceleration has a direction, and tabulating it is not cheap}
\label{sec:controller-direction}

A tail criterion predicts $\lVert\Delta\mathbf a\rVert$. The coupling term of
Eq.~\eqref{eq:coordinate-update} needs $\Delta\mathbf a$ itself. A controller
supplied only with the magnitude can reach rung 3 of Table~\ref{tab:ladder} and
no further, by construction.

The direction varies on the model's own resolution scale. It is set by the
leading omitted harmonics, whose half-wavelength at the spacecraft's radius is
$\pi r/N$: about \ph{value}~km at $N=120$ and \ph{value}~km at $N=300$. A tabulation,
interpolation or fitted surrogate over $(r,\phi,\lambda)$ that does not itself
resolve that scale therefore cannot represent the direction, and one that does
resolve it carries a number of degrees of freedom comparable to the omitted
coefficients themselves, with the storage and evaluation cost that implies.
This is a statement about resolution and cost, not an impossibility claim about
any particular function class; Section~\ref{sec:ablations} measures it directly
with a surrogate built for the purpose and reports what it retains and what it
costs. \wpref{WP5}

\subsection{Band probe}

What is affordable is to evaluate the first few omitted bands and use their sum
as the direction:

\begin{equation}
  \hat{\mathbf v}(t,N)
  \;=\; -\,\gamma(N,r)\,\sum_{n=N+1}^{N+k}\mathbf a_n(t) ,
  \label{eq:probe}
\end{equation}

with $\gamma$ the ratio of the full omitted tail to the probed bands. Direction
comes from the probe; amplitude from $\gamma$.

$\gamma$ is taken from the model's \emph{measured degree-variance spectrum},
not from a fitted power law:

\begin{equation}
  \gamma(N,r) \;=\;
  \Biggl[\frac{\sum_{n>N}\sigma_a^2(n;r)}
              {\sum_{n=N+1}^{N+k}\sigma_a^2(n;r)}\Biggr]^{1/2} ,
  \label{eq:gamma}
\end{equation}

with $\sigma_a(n;r)$ the per-degree acceleration RMS built from the coefficient
degree variances $P_n$. The distinction matters because the lunar spectrum is
not a single power law --- the previous paper measured a spectral slope of
$2.13$ against an effective tail exponent of $1.76$ for the same field, and the
gap is a property of the spectrum's shape --- so a power-law extrapolation of
the tail carries a systematic error that a table of $P_n$ does not. And it
costs nothing in deployability: $P_n$ is roughly \ph{n} numbers, a
one-dimensional table.

That split is worth naming, because it is the paper's information statement in
its sharpest form. \textbf{What is cheap to carry is the spectrum; what is
expensive is the texture.} The amplitude of the omitted tail is a
one-dimensional function of degree and can simply be uploaded. Its direction is
a field on the sphere at resolution $\pi r/N$ and cannot; it has to be probed.

\paragraph{Cost, and why the cheap-looking figure is the wrong one.} At a point
where the field is \emph{already} being evaluated to degree $N$, extending to
$N{+}k$ costs $(N{+}k)^2-N^2\approx 2kN$ against $N^2$, that is $2k/N$ --- about
5\% at $N=120$, $k=3$ --- because the associated Legendre recursion is
incremental in degree and the probe inherits it.

That figure does not apply here, and saying so is the honest form of the cost
accounting. A look-ahead probe is taken at a point the integrator has \emph{not}
evaluated, and the recursion for $P_{nm}$ at a new location must run from
degree zero: obtaining band $n$ alone is not an $\mathcal{O}(n)$ operation in
any standard synthesis. \textbf{Each look-ahead probe point therefore costs
approximately one full synthesis}, not an increment on one.

Two properties of the construction do reduce the cost, but both act
\emph{within} a probe point rather than between them.

First, the band stack is \emph{shared} across candidates. Evaluating the individual bands
$\mathbf a_n$ once for $n$ up to $N_{\max}{+}k$ yields
$\hat{\mathbf v}(t,N)$ for \emph{every} candidate $N\le N_{\max}$ by partial
summation, at no further field evaluation. The candidates are therefore not
independent probes but partial sums of one stack.

Second, the candidate set is a window. The offline plan supplies a nominal
degree for each decision interval; writing $j(q)$ for its position in the
ordered candidate grid $\mathcal{N}=\{N_1<N_2<\dots\}$, the online decision
searches only

\begin{equation}
  \mathcal{N}_q \;=\; \{\,N_{j(q)-\delta},\,\dots,\,N_{j(q)+\delta}\,\} ,
  \label{eq:window}
\end{equation}

so that $\delta$ is a half-width in \emph{candidate-grid indices}, not in
degrees. Fixing it that way keeps the controller's complexity ---
$2\delta+1$ candidates per decision --- invariant when the grid is refined or
coarsened; the induced span in degrees follows from the grid and is reported
alongside the cost, since it is the span and not the index count that sets the
band stack. The stack runs from $N_{j-\delta}{+}1$ to $N_{j+\delta}{+}k$, so
the incremental cost is $\approx 2N(\Delta_{\mathrm{span}}+k)$ with
$\Delta_{\mathrm{span}}=N_{j+\delta}-N_{j-\delta}$, rather than $2kN$.

This is the plan/probe split made concrete: the plan sets the coarse
allocation, which needs the budget and the sensitivity and is smooth; the probe
corrects within a window, which needs the local texture. The window is a
declared parameter, swept in Section~\ref{sec:ablations}, and it handicaps the
controller relative to the benchmark --- which searches the full $\mathcal{N}$
--- so the cost of the restriction appears in the capture fraction rather than
being hidden.

\paragraph{What the overhead actually is.} Covering a decision interval takes
$n_{\mathrm{probe}}\approx\Delta t_{\mathrm{dec}}/\tau_{\mathrm{corr}}$ probe
points, and the interval contains $n_{\mathrm{RHS}}\approx
\Delta t_{\mathrm{dec}}\,r_{\mathrm{RHS}}$ right-hand-side calls, so the
overhead is

\begin{equation}
  \frac{n_{\mathrm{probe}}}{n_{\mathrm{RHS}}}
  \;\approx\; \frac{1}{\tau_{\mathrm{corr}}\;r_{\mathrm{RHS}}}
  \;=\; \frac{N v}{\pi r\, r_{\mathrm{RHS}}} ,
  \label{eq:probe-overhead}
\end{equation}

which is \emph{independent of the decision cadence}: probing twice as often
over intervals twice as short buys nothing.

Equation~\eqref{eq:probe-overhead} is a local rate, and the number that belongs
in a budget is its arc integral,

\begin{equation}
  \frac{B_{\mathrm{probe}}}{B_{\mathrm{tot}}}
  \;=\; \frac{1}{N_{\mathrm{RHS}}}\int_0^T \frac{\dd t}{\tau_{\mathrm{corr}}(t)}
  \;=\; \frac{1}{N_{\mathrm{RHS}}}\int_0^T
        \frac{N(t)\,v(t)}{\pi\,r(t)}\,\dd t ,
  \label{eq:probe-overhead-arc}
\end{equation}

and the two are not interchangeable. Where $\tau_{\mathrm{corr}}$ is shortest
--- at perilune, where the degree is highest --- the vehicle spends the least
time, and the integrator's call rate is highest there as well, so both factors
move together. Reading the local rate at perilune against an arc-averaged call
rate overstates the cost substantially.

Evaluated on the archived reference arcs, Eq.~\eqref{eq:probe-overhead-arc}
gives \ph{pct} for a near-circular low orbit and \ph{pct} for the eccentric
geometries of Table~\ref{tab:populations}, the eccentric cases being
\emph{cheaper} because their long apolune stretches need almost no probing.
Those estimates hold $r_{\mathrm{RHS}}$ fixed at its arc average and therefore
still err high; the measured values are in
Table~\ref{tab:probe-cost}. The probe remains the dominant term in $B_+$
--- against \ph{pct} for the pilot arc and \ph{pct} for switching --- which is
why it is reported as a curve rather than absorbed.

Two consequences follow, and neither is a footnote. The overhead is reported as
a \emph{measured curve} rather than a number: Table~\ref{tab:probe-cost} gives
the effective direction accuracy $\kappa_{\mathrm{eff}}$ obtained by probing at
$n_{\mathrm{probe}}$ points, averaged over the interval, against the cost that
buys it. And a cheap endpoint of that curve is declared in advance:
\Clite\ raises the degree of the interval's \emph{first} right-hand-side call
--- a point the integrator evaluates anyway, so $2k/N$ genuinely applies ---
and accepts that the direction is known only at the interval's start, for
$\tau_{\mathrm{corr}}$ of its length. \wpref{WP5, WP16}

\paragraph{Switching is not the expensive part.} A decision boundary makes the
right-hand side discontinuous, and one might expect that to dominate. It does
not. What the previous campaign measured is the aggregate: a radially binned
schedule cost \ph{pct} extra right-hand-side calls over seven days. Converting
that to a per-switch figure requires counting the bin crossings, which we do
from the archived degree histories rather than estimating; on \ph{n} crossings
it is \ph{value} extra calls per switch. At
$\Delta t_{\mathrm{dec}}=\ph{value}$~s there are \ph{n} boundaries per arc, so
the switching term is \ph{pct} --- roughly an order of magnitude below the
probe.

\begin{table}[!htbp]
  \centering\small
  \caption{The cost--accuracy curve of the online probe, measured rather than
  counted. $\kappa_{\mathrm{eff}}$ is the direction accuracy averaged over the
  decision interval when the interval is probed at $n_{\mathrm{probe}}$ points;
  the cost column is the share of $B_{\mathrm{tot}}$ the probe consumes, which
  Eq.~\eqref{eq:matching} subtracts from the controller's gravity budget. The
  first row is \Clite, whose single probe is co-located with a right-hand-side
  call and therefore costs an increment rather than a synthesis.
  \ph{fill from WP5, WP16}}
  \label{tab:probe-cost}
  \begin{tabular}{@{}lccccc@{}}
    \toprule
    Variant & $n_{\mathrm{probe}}$ & $\delta$ (idx) & $\Delta_{\mathrm{span}}$ (deg)
      & $\kappa_{\mathrm{eff}}$ & cost, share of $B_{\mathrm{tot}}$ \\
    \midrule
    \Clite & 1 (co-located) & \ph{n} & \ph{n} & \ph{value} & \ph{pct} \\
    \Cplan & \ph{n} & \ph{n} & \ph{n} & \ph{value} & \ph{pct} \\
    \Cplan & \ph{n} & \ph{n} & \ph{n} & \ph{value} & \ph{pct} \\
    \bottomrule
  \end{tabular}
\end{table}

\begin{figure}[!tbp]
  \centering
  \ph{figures/fig\_oa\_probe.pdf}
  \caption{Probe direction accuracy $\kappa=\cos\angle(\hat{\mathbf v},
  \Delta\mathbf a)$ against probe depth $k$, truncation degree $N$ and
  altitude. The dashed line is the level below which the coupling term cannot
  be carried (Section~\ref{sec:controller-limits}).}
  \label{fig:probe}
\end{figure}

\ph{one paragraph} Median $\kappa$ is \ph{value} at $k=1$ and \ph{value} at
$k=3$; it \ph{degrades/holds} with altitude and \ph{degrades/holds} with $N$.
The refresh interval required to keep $\kappa$ above \ph{value} is
\ph{value}~s, consistent with Eq.~\eqref{eq:decorrelation}.

\paragraph{What $\kappa$ does and does not measure.} It is tempting to read
$\kappa$ as the fraction of the coupling term the probe retains. It is not,
and the difference matters enough to state. The term is
$2\,\mathbf u_i(N)^{\top}\mathbf c_i$ with
$\mathbf u_i=\Phi\mathbf B\,\Delta\mathbf a\,\Delta t_i$, so writing
$\mathbf z=(\Phi\mathbf B)^{\top}\mathbf c_i$ the retained fraction is
$\langle\hat{\mathbf v},\mathbf z\rangle/
 \langle\Delta\mathbf a,\mathbf z\rangle$: it depends on the angles the true
and estimated defects make with $\mathbf z$, not on the angle between them,
because $\Phi\mathbf B$ does not preserve angles. Where
$\Delta\mathbf a\perp\mathbf z$ the true contribution vanishes and any probe
error produces a spurious one.

$\kappa=1$ still forces the retained fraction to one, and a low $\kappa$ still
destroys it, so $\kappa$ is a leading indicator rather than the quantity
itself --- and it has the practical advantage of being measurable without a
schedule and without propagating. We therefore report both: $\kappa$ as the
leading indicator, and the retained fraction, computed with $\mathbf c_i$ from
the \Asens\ schedule, alongside it. \emph{Both are diagnostics.} Neither
carries a stopping rule; the deployability decision is the capture fraction of
Section~\ref{sec:campaign}, which measures the thing itself rather than a
proxy for it. \wpref{WP5}

\subsection{Planning: forward--backward allocation}
\label{sec:controller-ifbda}

The coupling term $\mathbf c_i$ of Eq.~\eqref{eq:cross-term} depends on the
degrees chosen at \emph{every other} epoch. A single backward sweep of an
adjoint cannot supply it, because there is no schedule yet to sweep against.
The controller therefore iterates.

\begin{algorithm}[!htbp]
\caption{Forward--backward degree allocation on a pilot arc}
\label{alg:ifbda}
\begin{algorithmic}[1]
\Require pilot arc; low-degree $\Phi$; tail model; budget $B$
\State $N^{(0)} \gets$ separable sensitivity-weighted schedule (\Asens\ form)
\For{$j = 1 \dots J$}
  \State \textbf{forward:} accumulate $\hat{\mathbf u}_i$ from
    $\hat{\mathbf v}(t_i,N^{(j-1)}_i)$; form prefix sums
  \State \textbf{backward:} form the suffix blocks $\mathbf A_i$ and the cross
    terms $\mathbf c_i$
  \State \textbf{update:} $N^{(j)}$ by Eq.~\eqref{eq:coordinate-update};
    bisect $\lambda$ to the budget
\EndFor
\State \Return cancellation target $\mathbf c$, diagonal kernel $\mathbf K_i$,
  multiplier $\lambda$
\end{algorithmic}
\end{algorithm}

$J$ is a free parameter and is swept in Section~\ref{sec:ablations};
\ph{value} iterations suffice on \ph{n} of \ph{n} panel orbits.

\subsection{Probing forward, not backward}
\label{sec:controller-causality}

A decision interval is $\Delta t_{\mathrm{dec}}\approx\ph{value}$~s long and the
omitted direction decorrelates in $\Delta t_{\mathrm{probe}}\approx\ph{value}$~s.
A rule that accumulated probes taken \emph{inside} an interval and applied them
at the \emph{next} boundary would therefore select a degree from a direction
that has already decorrelated --- exactly the failure the resolution-scale
argument of Section~\ref{sec:controller-direction} predicts. The controller
does the opposite: at each boundary it probes \emph{ahead}, over the interval it
is about to commit to.

That is legitimate, and the reason is that two different coherence questions are
involved and they have different answers.

\begin{table}[!htbp]
  \centering\small
  \caption{Two coherence scales. The long-horizon failure is what forces the
  plan/probe split of Section~\ref{sec:controller-split}; the short-horizon
  success is what makes forward probing admissible.}
  \label{tab:coherence}
  \begin{tabular}{@{}p{0.24\textwidth}p{0.30\textwidth}p{0.36\textwidth}@{}}
    \toprule
    Horizon & State prediction error & Against $\pi r/N$ \\
    \midrule
    One decision interval,
    $\Delta t_{\mathrm{dec}}$ & the vehicle's own short-arc prediction,
      \ph{value}~m & far below; forward probing valid \\
    \addlinespace
    Whole arc, pilot against
    flown & \ph{value}~km after $T_{\mathrm{coh}}$ & far above; a pre-computed
      $\hat{\mathbf v}$ is invalid \\
    \bottomrule
  \end{tabular}
\end{table}

\paragraph{What predicts the probe locations.} At $t_q$ the degree for the
upcoming interval has not been chosen, so the integrator has not yet advanced
through it, and the probe points have to come from somewhere else. Two facts
make this a small requirement rather than a circularity.

First, the probe locations barely depend on the degree at all. Two candidate
degrees differ in acceleration by the difference of their truncation defects,
so over one interval they separate by at most
$\tfrac12\lVert\Delta\mathbf a\rVert\Delta t_{\mathrm{dec}}^2$, which for
$\lVert\Delta\mathbf a\rVert\sim\ph{value}$ and
$\Delta t_{\mathrm{dec}}=\ph{value}$~s is \ph{value}~mm. Against $\pi r/N$ that
is \ph{value} orders of magnitude below the scale on which the direction
changes. The predictor therefore does not need to know $N_q$, which is what
breaks the circle.

Second, what the predictor must achieve is a position accurate to well inside
$\pi r/N$ --- kilometres --- not an accurate trajectory. The controller uses a
two-body propagation of the current state over the interval, which requires no
field evaluation and is analytic. Its error is dominated by the perturbing
acceleration it omits, $\tfrac12\lVert\mathbf a_{\mathrm{pert}}\rVert
\Delta t_{\mathrm{dec}}^2\approx\ph{value}$~m, again far inside the tolerance.
Its cost is \ph{value} and enters $B_+$ with the rest.

\code{abl-predictor} replaces it with a low-degree micro-propagation at
$N_{\mathrm{plan}}$ and reports the change in $\kappa$ and in the final error,
so the choice is measured rather than assumed \wpref{WP5}.

The controller therefore evaluates the shared band stack at
$n_{\mathrm{probe}}=\lceil\Delta t_{\mathrm{dec}}/\Delta t_{\mathrm{probe}}\rceil$
points spanning the upcoming interval, along that predicted path, and sums
their transported contributions to obtain each candidate's
$\sum_{i\in I_q}\mathbf u_i(N)$. Nothing from before the boundary is reused for
the direction.

\begin{algorithm}[!htbp]
\caption{Online decision at boundary $t_q$}
\label{alg:online}
\begin{algorithmic}[1]
\Require plan $(\mathbf c,\mathbf K_i,\lambda,N_{\mathrm{plan}})$; window half-width
  $\delta$ (in candidate-grid indices); depth $k$; spent work $W$
\State predict the state at $n_{\mathrm{probe}}$ points spanning
  $[t_q,t_{q+1})$ by two-body propagation of the current state
  \Comment{no field evaluation; degree-independent}
\State evaluate bands $N_{\mathrm{plan}}(t_q)-\delta+1 \dots
  N_{\mathrm{plan}}(t_q)+\delta+k$ at those points \Comment{one shared stack}
\For{$N \in \mathcal{N}_q$}
  \State $\hat{\mathbf v}(t_i,N) \gets -\gamma(N,r_i)\sum_{n=N+1}^{N+k}
    \mathbf a_n(t_i)$ \Comment{partial sums, no new evaluation}
\EndFor
\State $N_q \gets$ Eq.~\eqref{eq:coordinate-update} over $\mathcal{N}_q$ with
  $\mathbf c$ from the plan
\State $\lambda \gets \lambda + \mathrm{gain}\cdot(W - W_{\mathrm{sched}}(t_q))$
  \Comment{budget feedback, Section~\ref{sec:controller-feedback}}
\State hold $N_q$ over $[t_q,t_{q+1})$
\end{algorithmic}
\end{algorithm}

Holding the degree across the interval also gives the controller exactly the
switching freedom the benchmark of Section~\ref{sec:problem} has, so the two
differ in information and not in resolution.

\dnote{If the decision grid has to shrink toward
$\Delta t_{\mathrm{probe}}$ for the coupling term to be worth anything, say so:
that is a finding about the controller's required cadence, and it should be
reported with the integrator cost it provokes rather than hidden by choosing
$\Delta t_{\mathrm{dec}}$ after the fact. The sweep in
Section~\ref{sec:ablations} is what decides this.}

\subsection{Why the plan and the probe are split}
\label{sec:controller-split}

Algorithm~\ref{alg:ifbda} runs on a pilot arc, which is not the arc that will
be flown. Whether that matters depends on which quantity is being carried
across, and the two behave differently.

The local omitted vector $\hat{\mathbf v}$ is field texture and decorrelates
over $\pi r/N$. A pilot arc's position error exceeds that scale after
$T_{\mathrm{coh}}=\ph{value}$, measured in Section~\ref{sec:controller-limits},
which is \ph{fraction} of the arc --- the long-horizon row of
Table~\ref{tab:coherence}. A schedule that fixed the probe at plan time would
therefore be using a direction that is no longer the right one, and the
controller does not do that: $\hat{\mathbf v}$ is evaluated online, at the
vehicle's own state, one interval at a time
(Section~\ref{sec:controller-causality}).

The cancellation target $\mathbf c$ is an accumulated quantity. The previous
paper measured the defect to be coherent rather than incoherent --- growth
exponents of \ph{range} against the $t^{3/2}$ of a zero-mean impulse null ---
so the accumulated displacement is dominated by a sustained trend rather than
by local texture, and it is stable between the pilot and the reference arc to
\ph{value} \wpref{WP6}. The same is true of the sensitivity weight and of the
multiplier.

The controller is accordingly split:

\begin{center}
\begin{tabular}{@{}ll@{}}
  \toprule
  Computed offline, on the pilot arc & $\mathbf c$, $\mathbf K_i$, $\lambda_0$,
    $N_{\mathrm{plan}}$, \emph{indexed by orbital phase} \\
  Computed online, ahead of each boundary & $\hat{\mathbf v}$ by
    Eq.~\eqref{eq:probe}, hence $N_q$ within $\pm\delta$ of the plan \\
  \bottomrule
\end{tabular}
\end{center}

and neither half reads the reference field or the reference arc.

\paragraph{The plan is indexed by phase, not by time.} This is forced by the
same coherence argument, applied to the smooth quantities rather than to the
texture. A loose $N=40$ pilot arc drifts along track by of order \ph{value}~km
over seven days, which at low-lunar-orbit speed is a \emph{timing} offset of
\ph{value}~s. A perilune passage lasts a few minutes. Since $\mathbf c$ and
$\mathbf K_i$ are dominated by exactly those passages, a plan indexed by
absolute time would
apply the perilune cancellation target before or after the perilune it was
computed for --- an error of order the thing being corrected.

Indexing by (revolution number, phase within revolution) removes the
sensitivity: the vehicle knows its own phase, and a \ph{value}~s offset is
\ph{value} of a revolution. It is not a return to a single-valued function of
radius, because the revolution index carries the remaining horizon; that is
what distinguishes the plan from the radial rules of
Table~\ref{tab:ladder}. \wpref{WP6, WP9}

\subsection{Budget feedback}
\label{sec:controller-feedback}

A per-call budget does not survive contact with an adaptive integrator: the
solver takes short steps where a concentrated schedule runs its highest
degrees, so nominal equality becomes a realized overspend --- a median
\ph{pct} in the previous campaign, and a factor of \ph{value} at half the
budget. \Cfb\ carries the spent work as a state, accumulating $N^2$ over the
right-hand-side calls actually made and adjusting $\lambda$ so that the
declared budget is met end to end rather than per call. \ph{one paragraph on
the update law and its gain.} The realized-to-nominal ratio it achieves is
\ph{value} against \ph{value} for the open-loop schedule.

\subsection{Where the controller is not the benchmark}
\label{sec:controller-limits}

Four separations are measured rather than argued, and each is an ablation in
Section~\ref{sec:ablations}: the tail amplitude in place of the true defect
magnitude; the probe direction in place of the true direction; the pilot arc in
place of the reference arc; and the low-degree $\Phi$ in place of the exact
one. \ph{one paragraph summarizing which dominates.}

\dnote{This is where a low $\kappa$ would first show. It is a warning and not a
stopping rule --- $\kappa$ is not the retained fraction, so gating on it would
risk an early false negative --- but if the capture fraction then confirms it,
the coupling term cannot be carried, \Cplan\ collapses onto rung 3, and the
paper's constructive claim reduces to a sensitivity-weighted controller ---
which is still new relative to the previous work, but is a different and
smaller claim. Draft that version of this section before the number arrives.}
